% =====================================================================
% CAPITULO 6 - CONCLUSIONES
% =====================================================================
% 6.1 cierra los objetivos O1-O6 de la seccion 1.3.2 uno a uno. Si alli se
% toca un objetivo, hay que tocar su cierre aqui.
%
% ESTILO (31-ago): sin negritas en el texto corrido, sin incisos entre rayas
% salvo cuando la frase los pida, y sin florituras.
% =====================================================================

\section{Cumplimiento de los objetivos}
\label{sec:objetivos-cumplidos}

Los seis objetivos de la sección~\ref{sec:objetivos} se cierran a continuación
con la evidencia que los respalda. Dos se cumplieron de forma distinta a la
prevista y se indica cómo.

\paragraph{O1. Pipeline reproducible.} Cumplido. Siete días procesados desde la
captura en bruto hasta el conjunto etiquetado, con la conexión TCP como unidad y
cinco representaciones por conexión (sección~\ref{sec:pipeline}). El
almacenamiento es autocontenido, lo que permitió derivar dos vistas nuevas, la
ráfaga y el inter-arribo, sin volver a procesar unas capturas que ya se habían
borrado por falta de espacio en disco.

\paragraph{O2. Ground-truth verificada.} Cumplido, y con hallazgo. La etiqueta
de los siete días se reconstruyó sobre las capturas y se contrastó con la
documentación oficial, que resultó incorrecta en tres de ellos: la dirección del
atacante SSH del 14-02, el puerto del ataque lento del 16-02 y la continuidad de
la actividad de la botnet del 02-03 (sección~\ref{sec:ground-truth}).

\paragraph{O3. Comparación con protocolo único.} Cumplido, con el resultado
opuesto al esperado. Las tres representaciones se compararon con la misma
métrica, el mismo balanceo y el mismo tratamiento del benigno sobre los seis
días. El objetivo pedía establecer si la comparación discrimina, y la respuesta
es que no lo hace: ninguna vista baja de $0.9186$ y el análisis de contenido no
baja de $0.9734$ (sección~\ref{sec:saturacion},
tabla~\ref{tab:saturacion}). El objetivo se cumple; la hipótesis que lo motivaba,
no.

\paragraph{O4. Por qué acierta cada representación.} Cumplido. La saliencia
localiza el 31.1\,\% de la importancia en los primeros 48 bytes, que es la zona
del \emph{banner} en claro (sección~\ref{sec:saliency}), y la prueba de evasión
lo confirma hundiendo el \emph{recall} del contenido de $1.0000$ a $0.0025$
en el día de la botnet, con la excepción del régimen en claro
(sección~\ref{sec:evasion}, tabla~\ref{tab:evasion-cap5}).

\paragraph{O5. Generalización.} Cumplido, y con corrección de rumbo. Se midió
entre los seis pares de días volumétricos y entre los dos conjuntos de datos. El
componente que falla está identificado, y no es la representación sino el
umbral: el AUC no baja de $0.9724$ mientras el macro-F1 con umbral fijo da
$0.3333$ (sección~\ref{sec:cross-dia}, tabla~\ref{tab:calibracion}). Se aporta
además la receta correspondiente, que son 25 etiquetas del dominio nuevo.

\paragraph{O6. Contraste con el estado del arte.} Cumplido. Se implementaron un
\emph{transformer} sobre bytes y una red de atención cruzada, y se compararon a
presupuesto de parámetros comparable. Ninguna mejora los resultados
(sección~\ref{sec:sota-contraste}). El límite no está en la capacidad del
modelo.

\section{Conclusiones principales}
\label{sec:conclusiones-principales}

\paragraph{La evaluación dentro del día está saturada y no discrimina.} Es la
conclusión que reordenó el trabajo. Sobre estos conjuntos de datos, cualquiera
de las representaciones detecta cualquiera de los ataques, así que un
macro-F1 alto obtenido así no mide capacidad de detección. La causa que primero
se propuso, la memorización del host atacante, se contrastó y resultó falsa: lo
que las vistas aprenden es la herramienta con que se generó el ataque.

\paragraph{La detección robusta sí exige combinar contenido y conducta, pero
quien lo demuestra no es la comparación de aciertos.} Lo demuestran la evasión,
que hunde el contenido salvo donde la firma es genuina, la transferencia entre
dominios y los ataques que no dejan contenido que analizar. La tesis inicial del
trabajo sobrevive; la evidencia que la sostiene es otra.

\paragraph{El éxito del análisis de contenido suele ser una huella de la
herramienta, no del ataque.} Un atacante lo anula cambiando de biblioteca, sin
modificar su técnica. La excepción es el régimen en claro, donde los bytes que
el modelo usa son el ataque, y por eso resiste.

\paragraph{Hay ataques estructuralmente invisibles al análisis de contenido.} No
son ataques difíciles, son invisibles: el 94.4\,\% de las conexiones del
reconocimiento de la Infiltration no lleva un solo byte de aplicación. Ninguna
mejora del modelo lo arregla, porque no hay datos que mejorar. Es el argumento
más fuerte a favor de combinar vistas y el único que no depende de comparar
métricas.

\paragraph{Entre dominios lo que se rompe primero es la calibración.} Y se
arregla con 25 etiquetas, sin reentrenar. Es la conclusión con aplicación
práctica más directa de toda la memoria.

\paragraph{Sin ataques etiquetados, el rendimiento cae entre veinte y cincuenta
puntos.} Todas las cifras supervisadas de esta memoria suponen disponer de
ejemplos etiquetados del mismo tipo de ataque. Al retirar esa suposición y
entrenar solo con tráfico benigno, el área bajo la curva ROC baja a $0.6946$
en el régimen cifrado y a $0.5402$ en el volumétrico
(sección~\ref{sec:anomalia}). Esa es la distancia real entre lo que se mide sobre
un conjunto etiquetado y lo que un detector encontraría ante una amenaza nueva.

\paragraph{El enfoque híbrido aporta por dónde se une, no por disponer de las dos
fuentes.} Concatenar las entradas crudas da un resultado peor que la mejor de las
dos ramas por separado, porque las 257 dimensiones del contenido diluyen los
cuatro escalares de flujo. La fusión tardía lo evita
(sección~\ref{sec:fusion}).

\paragraph{El cuello de botella está en el dato y en el protocolo, no en el
modelo.} Ninguna arquitectura avanzada mejoró nada, y ahora se sabe por qué: no
quedaba margen que recoger.

\section{Lecciones metodológicas}
\label{sec:lecciones}

Esta sección recoge los errores que el trabajo cometió y corrigió. Se incluye
porque cada uno produjo un resultado plausible antes de ser detectado, y esa es
la clase de fallo que una memoria no debería ocultar.

\paragraph{El tráfico benigno contaminado, dos veces.} En el día del SSH, 1.589
de los 2.022 flujos etiquetados como benignos eran del propio atacante; en el
del DoS, el 65.1\,\% del benigno HTTP. Se estaba comparando el ataque consigo
mismo. El error produjo un techo de $0.78$ que durante semanas se interpretó
como una propiedad del problema. Un número bajo no levanta sospechas del mismo
modo que un $1.0000$, y ahí está la trampa.

\paragraph{Una ley inducida de dos puntos.} El comportamiento de la
generalización entre días se enunció con dos días, resultó falso al añadir el
tercero y volvió a cambiar al medirlo con AUC en lugar de con F1. Tres
versiones, dos de ellas equivocadas. La lección tiene dos partes: no inducir
reglas con dos observaciones, y no juzgar la transferencia con una métrica que
dependa de un umbral fijo.

\paragraph{Métricas que no eran lo que decían.} Las cifras del caso difícil se
citaron durante semanas como macro-F1 cuando eran exactitud en validación
cruzada, en un capítulo que dedica una sección a explicar por qué la exactitud
engaña. Y una figura comparaba dos conjuntos de características distintos bajo
la misma etiqueta, lo que exageraba un efecto real. Antes de citar una cifra hay
que comprobar qué métrica es y qué contiene.

\paragraph{Una intuición sobre las características que la ablación desmintió.}
Se había escrito que el coeficiente de variación local era la característica
clave del \emph{beaconing}, porque mide la regularidad de los intervalos y esa es
la intuición evidente. La ablación mostró que la característica cuya ausencia más
cuesta es la media local, que la mejor en solitario es el intervalo anterior en
crudo, y que una de las seis, \texttt{regularidad}, empeora el conjunto
(sección~\ref{sec:interarribo}). Medir la aportación de cada
característica antes de darla por buena, porque alguna puede restar.

\paragraph{Una hipótesis propia, contrastada y refutada.} La explicación que se
dio a la saturación, la memorización del host atacante, era razonable y era
falsa. Se diseñó el experimento que podía tumbarla, se verificó con un control
positivo que el experimento era sensible, y se tumbó
(sección~\ref{sec:por-que-satura}). Es el episodio del que más se aprendió.

\paragraph{Una consecuencia práctica.} Estas correcciones motivaron una
herramienta de verificación que contrasta cada cifra de la memoria contra la
fuente de resultados, incluida la procedencia de las tablas. No sustituye a
leer, pero impide que un número se quede desfasado en silencio.

\section{Limitaciones}
\label{sec:limitaciones}

\begin{itemize}
  \item Un solo conjunto de datos como eje. CIC-IDS2017 se usa como validación
        complementaria, no como segundo eje.
  \item Un único host atacante en cinco de los seis días. Es lo que impide
        contrastar la hipótesis del atajo por identidad de host más allá del día
        del DDoS. En los otros cinco queda sin contrastar.
  \item Las características de ráfaga cuentan conexiones por origen. Aunque el
        experimento de exclusión de atacantes muestra que generalizan a un host
        no visto, la reserva merece declararse.
  \item La evasión probada es de una sola clase, el camuflaje de los primeros
        bytes. Un atacante que espacie deliberadamente sus conexiones atacaría
        la vista conductual con el mismo esfuerzo, y ese caso no se ha medido.
        Ser invariante no equivale a ser robusto.
  \item La detección sin etiquetas de ataque se ensaya solo con dos métodos
        clásicos y sobre tres días. Es suficiente para acotar la distancia
        respecto al caso supervisado, pero no para concluir nada sobre el estado
        del arte en detección de anomalías.
  \item Se contrasta la arquitectura, no el preentrenamiento. La comparación con
        el \emph{transformer} no incluye el preentrenamiento masivo, que es
        buena parte de lo que aporta ET-BERT.
  \item No hay día homólogo de botnet, con lo que el resultado del día 02-03 no
        está validado contra el test cruzado que sí se aplicó al DoS.
  \item Los grafos y los modelos de espacio de estados quedan fuera de alcance,
        y los primeros son justamente los que capturarían el caso que este
        trabajo identifica como su límite.
\end{itemize}

\section{Trabajo futuro}
\label{sec:futuro}

Se ordenan por retorno esperado, justificando cada línea con un resultado de
esta memoria y no con su interés general.

\paragraph{Protocolo que rompa el atajo en todos los días.} El experimento de
exclusión de atacantes solo pudo aplicarse al día con diez orígenes. Capturas
con varios atacantes por escenario permitirían generalizarlo, y es la
continuación directa de la sección~\ref{sec:por-que-satura}.

\paragraph{Evasión sobre la vista conductual.} Un atacante que ralentice su
fuerza bruta es el equivalente conductual de camuflar el \emph{banner}. Hasta
medirlo, la robustez de la ráfaga está afirmada solo frente a una clase de
evasión.

\paragraph{Redes sobre grafos para el movimiento lateral.} Lo pide el resultado
de la sección~\ref{sec:limite}: el reconocimiento de la Infiltration es un
patrón topológico, un origen contra 612 destinos y sin carga útil, que ninguna
de las vistas de este trabajo puede representar.

\paragraph{Recalibración continua en despliegue.} La
sección~\ref{sec:cross-dia} muestra que 25 etiquetas bastan para reajustar el
umbral en un dominio nuevo. Convertirlo en un procedimiento periódico es una
línea de ingeniería con resultado medible.

\paragraph{Validación entre conjuntos de datos como norma.} No como
complemento. Es la única de las validaciones aplicadas que rompe la fuga de
identificadores (sección~\ref{sec:fuga}), y por tanto la única que garantiza que
lo medido sea detección y no memorización del laboratorio.
